%% kapa: exam=1 round=31 year=2020 period=2 form=A subject=accounting subject_ko=회계학 seq=20 number=60
%% source: Q-Net 감정평가사 기출문제 (cst00309) artlSeq=5207219 file=extracted/1008/5207219/A형/A형/2020년도 제31회 감정평가사 1차 시험_2교시_A형.pdf
%% conversion: kordoc 4.24.51 -> tools/gen1.py (auto); review=proofread-sonnet; flags=none
\begin{question}{60}
다음은 (주)감평의 수익 관련 자료이다.
\begin{framedtext}
○ (주)감평은 20x1년 초 (주)한국에게 원가 ￦50,000의 상품을 판매하고 대금은 매년 말 ￦40,000씩 총 3회에 걸쳐 현금을 수취하기로 하였다. \par
○ (주)감평은 20x1년 12월 1일 (주)대한에게 원가 ￦50,000의 상품을 ￦120,000에 현금 판매하였다. 판매계약에는 20x2년 1월 31일 이전에 (주)대한이 요구할 경우 (주)감평이 판매한 제품을 ￦125,000에 재매입해야 하는 풋옵션이 포함된다. 20x1년 12월 1일에 (주)감평은 재매입일 기준 제품의 예상 시장가치는 ￦125,000 미만이며, 풋옵션이 행사될 유인은 유의적일 것으로 판단하였으나, 20x2년 1월 31일까지 풋옵션은 행사되지 않은 채 소멸하였다. \par
\end{framedtext}
(주)감평이 20x2년에 인식해야 할 총수익은? (단, 20x1년 초 (주)한국의 신용특성을 반영한 이자율은 5\%이고, 계산금액은 소수점 첫째자리에서 반올림하며, 단수차이로 인한 오차가 있으면 가장 근사치를 선택한다.)
\par{\small\centering\begin{tabular}{|c|c|c|}\hline
기간 & \makecell{단일금액 ￦1의 현재가치\\(할인율=5\%)} & \makecell{정상연금 ￦1의 현재가치\\(할인율=5\%)} \\ \hline
3 & 0.8638 & 2.7232 \\ \hline
\end{tabular}\par}
\begin{choices}
  \item ￦0
  \item ￦120,000
  \item ￦125,000
  \item ￦128,719
  \item ￦130,718
\end{choices}
\end{question}
\answer{4}
